\author{Schmid Lothar}
\documentclass[11pt]{mybibreport}

\graphicspath{{../../../assets/images/}}


% Absatzformatierung
% \setlength{\parindent}{0pt}
% \setlength{\parskip}{0.6em}
% \renewcommand{\arraystretch}{1.8}
\begin{document}

\tableofcontents
\chapter{Allgemeines zum Buch}
\section{Author}
Der Schreiber dieses Briefes wird in Vers 1 als Simon Petrus genannt. Obwohl dieser Name direkt als Verfasser im Text steht, wurde dies nicht akzeptiert. Weil der Wortschatz und die Grammatik in diesem Brief anders war als die im 1. Petrusbrief, wurde von speziallisten behauptet, dass es zwei unterschiedliche Autoren sein muss.
\subsection{Verfassungszeit}
\subsection{Verfassungsort}


\chapter{Textanalyse}

\begin{bibelanalyse}{2. Petrus}{1}{1-11}
    \analysistexte{
       
    \versmarker{1}{
        \hauptsatz{\person[b, p]{Simon Petrus}, Knecht \konj[b, p]{und} Apostel \person[b, p]{Jesu Christi},}
        \nebensatz{an alle,} 
        \nebensatz{die mit uns denselben kostbaren Glauben \verbN[b, p]{empfangen haben}}
        \nebensatz[2]{\konj[b, p]{durch} die Gerechtigkeit unseres \person[b, p]{Gottes} \konj[b, p]{und} \person[b, p]{Heilands Jesus Christus}:}    
    }     
    \versmarker{2}{
        \nebensatz{\person[b, p]{Gott} \verbN[b, p]{gebe} euch viel Gnade und Frieden} 
        \nebensatz[2]{\konj[b, p]{durch} die Erkenntnis \person[b, p]{Gottes} \konj[b, p]{und} \person[b, p]{Jesu}, unseres \person[b, p]{Herrn}!}
    }
    
    \versmarker{3}{
        \hauptsatz{\konj[b, p]{Alles}, was zum Leben \konj[b, p]{und} zur Frömmigkeit \verbN[b, p]{dient}, }
        \nebensatz{\verbN[b, p]{hat} uns seine göttliche Kraft \verbN[b, p]{geschenkt}}
        \nebensatz[2] {\konj[b, p]{durch} die Erkenntnis dessen,}
        \nebensatz[3] {der uns \verbN[b, p]{berufen} \verbN[b, p]{hat}} 
        \nebensatz[2] {\konj[b, p]{durch} seine Herrlichkeit und Kraft.}
        }
    
    \versmarker{4}{
            \nebensatz[2] {Durch sie sind uns die kostbaren und allergrößten Verheißungen geschenkt,}
            \nebensatz[3] {damit ihr durch sie Anteil bekommt an der göttlichen Natur,}
            \nebensatz[2] {wenn ihr der Vergänglichkeit entflieht, die durch Begierde in der Welt ist.} 
        }
    \versmarker{5}{
            \hauptsatz{So wendet allen Fleiß daran}
            \nebensatz[2]{und erweist in eurem Glauben Tugend}
            \nebensatz[2]{und in der Tugend Erkenntnis}
        }
    \versmarker{6}{
            \nebensatz[2]{und in der Erkenntnis Mäßigkeit}
            \nebensatz[2]{ und in der Mäßigkeit Geduld}
            \nebensatz[2]{und in der Geduld Frömmigkeit}        
        }
    \versmarker{7}{
            \nebensatz[2]{und in der Frömmigkeit Brüderlichkeit} 
            \nebensatz[2]{und in der Brüderlichkeit die Liebe. }
        }
    \versmarker{8}{
            \hauptsatz{Denn wenn dies alles reichlich bei euch ist,}
            \nebensatz{wird's euch nicht faul und unfruchtbar sein lassen}
            \nebensatz{in der Erkenntnis unseres \person[b, p]{Herrn Jesus Christus}.}
        }
    \versmarker{9}{
            \hauptsatz{Wer dies aber nicht hat,}
            \nebensatz{der ist blind und tappt im Dunkeln und hat vergessen,}
            \nebensatz{dass er rein geworden ist von seinen früheren Sünden.} 
        }
    \versmarker{10} {
            \hauptsatz{Darum, Brüder und Schwestern,}
            \nebensatz{bemüht euch umso eifriger,}
            \nebensatz{eure Berufung und Erwählung festzumachen.}
            \nebensatz{Denn wenn ihr dies tut,}
            \nebensatz{werdet ihr niemals straucheln,}
        }
    \versmarker{11} {
            \nebensatz{und so wird euch reichlich gewährt werden}
            \nebensatz{der Eingang in das ewige Reich}
            \nebensatz{ unseres Herrn und Heilands Jesus Christus.} 
        }   
    }
    
    \begin{hinweise}
        
        \item \vers{1}{Petrus bezeichnet sich als Knecht, aber auch auch als ein Apostel also Botschafter Jesu Christi}
        \item \vers{1}{Sein Gruss geht an alle, die den gleichen Glauben wie Petrus haben. Am Ende des Briefes erwähnt er Paulus, darum spricht er hier in der Mehrzahl.}
        
    \end{hinweise}

    \begin{beobachtungen}
        
        \item \vers{1}{Da! Diese DA bezieht sich auf die vorgehenden Kapitel. }
        
    \end{beobachtungen}
    \begin{aussageabsicht}      
        1
        & 
        und
        & 
        Petrus bezeichnet sich als Knecht und Apostel
        \\
        \hline
       
    \end{aussageabsicht}
\end{bibelanalyse}

\end{document}